%% ---------------------------------------------------------------------------
%% Nushrat Naushin -- two-page resume (public)
%%
%% Build:  pdflatex resume.tex   (run twice so the page count is right)
%%
%% This is the SHORT, PUBLIC document: a two-page condensation of the master CV
%% that lives in the private Resume repository. It is deliberately trimmed --
%% no exhaustive coursework lists, no full publication list, no committee
%% rosters. When the master CV gains something significant, decide whether it
%% earns space here; most additions will not.
%%
%% Hard rule: this must stay at two pages. Check the page count after editing.
%% ---------------------------------------------------------------------------
\documentclass[10pt,letterpaper]{article}

\usepackage[T1]{fontenc}
\usepackage[utf8]{inputenc}
\usepackage{mathptmx}
\usepackage[margin=0.6in,top=0.55in,bottom=0.6in]{geometry}
\usepackage{enumitem}
\usepackage{titlesec}
\usepackage{xcolor}
\usepackage{microtype}

\definecolor{linknavy}{HTML}{14487F}
\definecolor{rulegray}{gray}{0.45}
\usepackage[colorlinks=true,urlcolor=linknavy,linkcolor=black,citecolor=black]{hyperref}

\emergencystretch=2.5em

\titleformat{\section}
  {\normalfont\normalsize\bfseries\raggedright}
  {}{0pt}{\MakeUppercase}[\vspace{-0.5em}\textcolor{rulegray}{\rule{\linewidth}{0.6pt}}]
\titlespacing*{\section}{0pt}{0.8em}{0.35em}

\setlist[itemize]{leftmargin=1.05em,topsep=0.1em,itemsep=0.06em,parsep=0pt,label=\textbullet}
\setlength{\parindent}{0pt}
\pagestyle{empty}

\newcommand{\entry}[3]{%
  \vspace{0.3em}%
  \noindent\textbf{#1}\hfill\textbf{\textit{#2}}\par
  \def\tempthird{#3}\ifx\tempthird\empty\else\noindent\textit{#3}\par\fi
  \vspace{0.05em}%
}

\newcommand{\skillline}[2]{\vspace{0.2em}\noindent\textbf{#1:} #2\par}

\begin{document}

%% ---------------------------------------------------------------- header ----
\begin{center}
  {\LARGE\bfseries NUSHRAT NAUSHIN}\\[0.35em]
  Albany, CA \quad$|$\quad \href{mailto:naushin264@gmail.com}{naushin264@gmail.com}
  \quad$|$\quad \href{https://sites.google.com/view/nushrat-naushin/}{website}
  \quad$|$\quad \href{https://www.linkedin.com/in/nushratn/}{LinkedIn}
\end{center}

%% --------------------------------------------------------------- summary ----
\section{Summary}

Materials scientist and metallurgical engineer with 4+ years in synchrotron X-ray spectroscopy and
microscopy, including two years as an Advanced Light Source Doctoral Fellow at Lawrence Berkeley
National Laboratory, plus 2.5 years of full-time undergraduate teaching and instructional
laboratory development. Designs and runs synchrotron experiments (XMCD, XMLD, XPEEM) on thin films
and complex alloys, builds statistically rigorous Python analysis pipelines for large experimental
datasets, and brings hands-on thin-film deposition, UHV maintenance, and classical mechanical and
metallurgical testing.

%% ------------------------------------------------------------ experience ----
\section{Experience}

\entry{Advanced Light Source, Lawrence Berkeley National Laboratory}{Sep 2024 -- Aug 2026}{ALS Doctoral Fellow}
\begin{itemize}
  \item Independently operate beamlines 4.0.2, 6.3.1.1, and 11.0.1.1; train new users on XMCD and
        XPEEM measurements.
  \item Lead 20+ beamtime experiments end to end: experimental planning, in-situ decisions during
        acquisition, and post-experiment data validation.
  \item Built XMCD sum-rule analysis with six-source Monte Carlo error propagation, and XPEEM
        asymmetry-map pipelines (drift correction, MCP noise removal, cross-registration, domain
        heterogeneity metrics) across temperature series.
  \item Developed and characterized reference multilayer thin films to improve beamline calibration
        and reproducibility.
  \item Co-authored 5+ successful beamtime proposals and contributed to multiple manuscripts.
\end{itemize}

\entry{University of California, Davis}{Sep 2021 -- Expected Sep 2026}{Graduate Student Researcher}
\begin{itemize}
  \item Deposit multi-principal element alloy thin films by DC magnetron sputtering and
        co-sputtering (AJA Orion 8); commission power supplies and maintain UHV systems, including
        turbomolecular pump repair.
  \item Lab superuser for the DC/RF sputtering system and Vibrating Sample Magnetometer; trained
        roughly 20 students on safe instrument operation.
  \item Correlate structure, composition, and magnetic behavior by combining magnetometry, XRD, and
        element-specific synchrotron spectroscopy.
  \item Lead laboratory safety programs, SOP development, and hazardous waste management; recipient
        of the College of Engineering Safety Excellence Award.
\end{itemize}

\entry{Khulna University of Engineering and Technology}{Jan 2019 -- Aug 2021}{Lecturer, Materials Science and Engineering}
\begin{itemize}
  \item Led outcome-based curriculum development and instruction for 8+ undergraduate theory and
        laboratory courses, typically about 60 students per class.
  \item Established and upgraded foundry, ceramics, electrochemical analysis, polymer and
        composites, and heat-treatment instructional laboratories, including equipment
        specifications and technical tender documents for major purchases.
  \item Mentored 10+ students through successful graduate school applications.
\end{itemize}

%% ------------------------------------------------------------- education ----
\section{Education}

\entry{University of California, Davis --- Ph.D., Materials Science and Engineering}{Sep 2021 -- Expected Sep 2026}{}
\begin{itemize}
  \item Element-specific magnetism and magnetic domain structure in multi-principal element alloy
        thin films; dissertation on composition-dependent magnetism in MnCoNi. Advanced to
        candidacy October 2023. PI: Prof.\ Roopali Kukreja.
\end{itemize}

\entry{Bangladesh University of Engineering and Technology --- B.S., Materials and Metallurgical Engineering}{Oct 2018}{}
\begin{itemize}
  \item Honours; GPA 3.79/4.00, ranked 6th of 52. Coursework across extractive and process
        metallurgy, casting, forming, and joining.
\end{itemize}

%% ---------------------------------------------------------------- skills ----
\section{Technical Skills}

\skillline{Thin-Film Synthesis}{DC/RF magnetron sputtering and co-sputtering (AJA Orion 8),
sol--gel synthesis; UHV operation, maintenance, and troubleshooting including turbomolecular pumps.}

\skillline{Synchrotron and Characterization}{XAS, XMCD, XMLD, XPEEM, XPCS, micro-/nano-diffraction;
XRR, XRD, VSM, SEM/EDS, AFM, MFM, TEM/electron diffraction.}

\skillline{Data Analysis}{Python (NumPy, Pandas, Jupyter, Matplotlib); Monte Carlo error
propagation, image registration and similarity metrics; LabVIEW, OriginPro, GenX, Gwyddion.}

\skillline{Mechanical and Metallurgical Testing}{UTM tensile testing; Rockwell and Brinell hardness;
Charpy and Izod impact; optical emission spectroscopy; metallography; failure analysis.}

\skillline{Tools}{Git, \LaTeX{}, AutoCAD, SolidWorks, Inkscape.}

%% ---------------------------------------------------------- publications ----
\section{Selected Publications}
\begin{itemize}
  \item N.\ Naushin \textit{et al.}, ``Element-specific magnetic moment and domain structure in
        magnetron-sputtered MnCoNi thin films,'' \textit{J.\ Magn.\ Magn.\ Mater.} (2026).
        \href{https://doi.org/10.1016/j.jmmm.2026.174250}{10.1016/j.jmmm.2026.174250}
  \item W.\ Beeson, N.\ Naushin \textit{et al.}, \textit{Adv.\ Funct.\ Mater.} (2025).
        \href{https://doi.org/10.1002/adfm.202424741}{10.1002/adfm.202424741}
  \item N.\ Z.\ Hagstrom, N.\ Naushin \textit{et al.}, \textit{Nature Communications} (submitted,
        2025).
  \item P.\ Sharma, N.\ Naushin, S.\ Rohila, and A.\ Tiwari, ``Magnesium Containing High Entropy
        Alloys,'' in \textit{Magnesium Alloys --- Structure and Properties} (2022).
        \href{https://doi.org/10.5772/intechopen.98557}{10.5772/intechopen.98557}
\end{itemize}

%% -------------------------------------------------------- awards and press ----
\section{Awards and Press}
\begin{itemize}
  \item ALS Doctoral Fellowship in Residence, Lawrence Berkeley National Laboratory (2024, 2025).
  \item UC Davis College of Engineering Graduate Excellence Award for Safety (2025).
  \item BUET Dean's Choice Award (2018); University Merit Scholarship (2015, 2017).
  \item LPS Qubit Collaboratory Quantum Computing Summer Short Course, Algorithms to Hardware
        (Jul--Aug 2025).
  \item Featured in
        ``\href{https://nobumichitamura.substack.com/p/chaos-by-design-how-mixing-metals}{Chaos by
        Design}'' by N.\ Tamura, ALS (2026);
        ``\href{https://mse.engineering.ucdavis.edu/news/college-honors-mse-graduate-students-excellence-safety}{College
        Honors MSE Graduate Students for Excellence in Safety},'' UC Davis (2025); and
        ``\href{https://mse.engineering.ucdavis.edu/news/international-womens-day-spotlight-uc-davis-women-materials-science-and-engineering}{International
        Women's Day Spotlight},'' UC Davis (2024).
\end{itemize}

%% ----------------------------------------------------------- conferences ----
\section{Conferences}
\begin{itemize}
  \item APS March Meeting --- Presenter (2025, 2026).
  \item ALS User Meeting --- Invited Speaker (2024); Poster Presenter (2026).
\end{itemize}

%% --------------------------------------------------------------- service ----
\section{Service and Mentorship}
\begin{itemize}
  \item Trained roughly 20 students on safe operation of the sputtering system and VSM as lab
        superuser, and mentored 14+ undergraduate and graduate students on advanced instrumentation.
  \item Mentored students from underrepresented backgrounds through the UC--HBCU and MESA
        initiatives, including intensive summer mentorship of a visiting undergraduate researcher.
  \item Advised 10+ former students on graduate school applications; all are now enrolled in
        graduate programs.
  \item Secretary, BUET Material Advantage Chapter: organized an inter-university materials poster
        symposium (50+ teams) and a STEM competition (500+ participants).
  \item Member, American Physical Society and The Minerals, Metals \& Materials Society.
\end{itemize}

\end{document}
